% Supplementary material (DAN-219, 2026-11-16, decision 7): moved verbatim out of paper/main.tex, app:production
% ("Circuit sizes, sufficiency to induce, training dynamics"), whose prose states the curves' statistics and points here.
% Not compiled with main.tex; \runfigure is defined in main.tex's preamble.

\begin{figure}[tbp]
\centering
% experiments/111-failure-anatomy a3_curves.py writes figures/run2026_training_curves.pdf.
\runfigure{training_curves}{circuit size per step, data loss per pass, members removed over the last 100 steps by layer; one row per model}
\caption{Production circuits shrink slowly by the 400-step budget, more at depth.
Training dynamics of the production fits by target depth band, one row per model (median and interquartile range).
Left: circuit size (latents above the keep threshold) per step.
Middle: the data loss (the Z, C, A, and R terms, with the rank-keep parts, penalty excluded), averaged over each pass through the training contexts.
Right: the share of members removed over the last 100 steps, by target layer (median and interquartile range).}
\label{fig:run2026-training-curves}
\end{figure}
